\usepackage{espresso}

% Значения перечислений (DoxyEnumFields) набираются таблицей longtable, а
% не longtblr из tabularray, как в espresso.sty (копии doxygen.sty из
% Doxygen). Doxygen набирает такую таблицу дважды, чтобы узнать её ширину,
% и на длинном перечислении (у esp::ps2::virtual_key около двухсот
% значений) tabularray 2023A из TeX Live 2023 исчерпывает память и
% регистры XeTeX. longtable набирает строки по мере чтения, без таких
% пределов, и переносит длинную таблицу на следующие страницы.
%
% Ширина колонок постоянная: имена — 35 % строки, описания — остальное.
% Первая строка — заголовок без вертикальных линий, как у Doxygen; на
% следующих страницах таблица начинается с линии. Строки тела Doxygen сам
% завершает \\ и \hline.
\newlength{\DoxyEnumNameWidth}
\newlength{\DoxyEnumTextWidth}
% #1 — колонки таблицы, #2 — их число, #3 — заголовок, #4 — тело.
\NewDocumentCommand{\DoxyEnumTable}{m m m +m}{%
  \par
  \begingroup
  \setlength{\LTleft}{0pt}%
  \setlength{\LTright}{\fill}%
  \begin{longtable}{#1}
  \multicolumn{#2}{l}{\tf{#3}} \\[1ex]
  \hline
  \endfirsthead
  \hline
  \endhead
  #4
  \end{longtable}%
  \endgroup
  \par\addvspace{6pt}%
}
\RenewDocumentEnvironment{DoxyEnumFields}{O{2} m +b}{%
  \setlength{\DoxyEnumNameWidth}{\dimexpr 0.35\linewidth - 2\tabcolsep\relax}%
  \ifnum#1>2
    % Имя, значение и описание; значение — по ширине содержимого.
    \setlength{\DoxyEnumTextWidth}{\dimexpr 0.45\linewidth - 2\tabcolsep\relax}%
    \DoxyEnumTable{|>{\raggedright\arraybackslash}p{\DoxyEnumNameWidth}|r|p{\DoxyEnumTextWidth}|}%
                  {3}{#2}{#3}%
  \else
    \setlength{\DoxyEnumTextWidth}{\dimexpr 0.65\linewidth - 2\tabcolsep - 3\arrayrulewidth\relax}%
    \DoxyEnumTable{|>{\raggedright\arraybackslash}p{\DoxyEnumNameWidth}|p{\DoxyEnumTextWidth}|}%
                  {2}{#2}{#3}%
  \fi
}{}

\PassOptionsToPackage{warn}{textcomp}
\usepackage{textcomp}
\usepackage[nointegrals]{wasysym}

% Packages required by doxygen
\makeatletter
\providecommand\IfFormatAtLeastTF{\@ifl@t@r\fmtversion}
% Suppress package identification of infwarerr as it contains the word "warning".
\let\@@protected@wlog\protected@wlog
\def\protected@wlog#1{\wlog{package info suppressed}}
\RequirePackage{infwarerr}
\let\protected@wlog\@@protected@wlog
\makeatother
\IfFormatAtLeastTF{2016/01/01}{}{\usepackage{fixltx2e}} % for \textsubscript
\ifPDFTeX
    \IfFormatAtLeastTF{2015/01/01}{\pdfsuppresswarningpagegroup=1}{}
\fi
